% !TeX root = ../main.tex
\section{Pixel detector}
\label{sec:pixel-detector}
The pixel detector comprises four barrel layers and ten discs at each end,
with \num{5666} modules and \num{12410} active chip patches. The inner two
barrels use single-chip modules and the outer two use quads. The endcaps combine
four inner single-module rings and two outer quad rings, with a larger active
aperture at the forward stations. Carbon sandwich staves and discs carry the
modules, cooling and local services. \Cref{fig:pixel-overview} shows the selected
layout and its connection to the outer service route.

\begin{figure}[!htbp]
\centering\includegraphics[width=\linewidth]{pixel-overview}
\caption{Pixel barrel and positive endcap, including local cooling, cable
routes and the barrel collection bay. The opposite end is reflected in $z$.}
\label{fig:pixel-overview}
\end{figure}

\subsection{Modules and readout}
\label{sec:pixel-modules}
Each hybrid module uses a planar n-in-p sensor, ITkPix v2 ASICs and a backside
flex. The \num{400} by \num{384} matrix has square
\qty{50}{\micro\metre} pixels, giving an active area of
\qty{20}{\milli\metre} by \qty{19.2}{\milli\metre} and
\num{153600} channels per chip. The approximate
\qty{20}{\milli\metre} by \qty{21}{\milli\metre} die reserves the remaining
space for chip periphery.\cite{rd53} A quad shares one sensor and flex between
four active patches, with \num{614400} channels. Guards and internal seams are
passive silicon. The substrate, active patches and occupied module body are
represented separately; their dimensions are given in
\cref{tab:pixel-sensor-dimensions}.

In n-in-p sensors the implants and guard structures face the readout chips;
the backside flex carries bias, power and signals. Edge insulation separates
the biased sensor from nearby chip surfaces.\cite{slim} The flex uses a
\qty{25}{\micro\metre} polyimide core and two
\qty{12.5}{\micro\metre} coverlays, represented together as
\qty{50}{\micro\metre} of polyimide. Two \qty{18}{\micro\metre} copper foils
at \qty{50}{\percent} coverage and \qty{25}{\micro\metre} epoxy complete
its effective material inventory. Sensor and ASIC silicon are each
\qty{150}{\micro\metre} thick, separated by a
\qty{20}{\micro\metre} bump standoff. Chip peripheries connect to exposed
flex pads through the selected \qty{25}{\micro\metre} aluminium bond-wire
concept and a module-side termination.

The stack runs from flex copper, polyimide, inner copper and epoxy through
sensor, bump standoff, ASIC and graphite contact toward the support.
The \qty{0.255}{\milli\metre} graphite shims connect the ASIC backs at
$w=\qty{0.245}{\milli\metre}$ to the body-back plane at
$w=\qty{0.5}{\milli\metre}$. Insulation separates the module's serial-power
potential from conductive supports. \Cref{fig:pixel-modules} combines the
module footprints and stack; the occupied body is \qty{1}{\milli\metre} deep.
Endcap single bodies are \qty{20.4}{\milli\metre} tangentially by
\qty{21.4}{\milli\metre} radially; quad bodies are
\qty{43.2}{\milli\metre} by \qty{44.2}{\milli\metre}, with ASIC peripheries
at the outer radial edges.

\begin{table}[!htbp]
\caption{Physical sensor and active-map dimensions in local $(u,v)$, in
\unit{\milli\metre}. Active bounding spans include the quad's internal seams;
active area is the sum of the separate patches.}
\label{tab:pixel-sensor-dimensions}
\centering\small
\begingroup\setlength{\tabcolsep}{3pt}
\begin{tabular}{@{}lrrrr@{}}
\toprule & Single barrel & Quad barrel & Single endcap & Quad endcap\\\midrule
Active bounding span & $19.2\times20.0$ & $38.6\times40.2$ & $20.0\times19.2$ & $40.2\times38.6$\\
Sensor substrate & $20.2\times20.4$ & $39.6\times40.6$ & $20.2\times19.4$ & $40.4\times38.8$\\
Guard per outer edge & $0.5,\ 0.2$ & $0.5,\ 0.2$ & $0.1,\ 0.1$ & $0.1,\ 0.1$\\
Active area [\unit{\milli\metre\squared}] & 384 & 1536 & 384 & 1536\\\bottomrule
\end{tabular}
\endgroup
\end{table}

\begin{figure}[!htbp]
\centering
\includegraphics[width=.87\linewidth]{pixel-modules}\par\smallskip
\includegraphics[width=.78\linewidth]{pixel-module-stack}\par\smallskip
\includegraphics[width=.92\linewidth]{pixel-module-edge}
\caption{Module footprints (top), the quad material stack through one chip
row (middle), and schematic bond access (bottom). Active patches, passive
silicon and occupied envelopes are distinct. Stack thicknesses use an
independent scale; the electrical detail illustrates discrete connections.}
\label{fig:pixel-modules}\label{fig:pixel-module-stack}
\end{figure}

\subsection{Barrel layout}
\label{sec:pixel-barrels}\label{sec:pixel-packing}
The four layer datums are at
\qtylist{31.5;57;96;182}{\milli\metre}. Single-chip staves tile the smaller
circumferences; quads reduce repeated module edges at larger radii. Adjacent
staves alternate radial levels, with the offsets and populations in
\cref{tab:pixel-barrels}. Their modules have no azimuthal tilt.
The \num{82} staves carry \num{3050} modules and \num{6794} chip patches.

Local $u$ is tangential, $v$ follows $z$, and $w$ points radially outward.
The \qty{20}{\milli\metre} active dimension follows $z$, placing the
approximately \qty{1.8}{\milli\metre} chip periphery at the transverse edge.
The longitudinal guard and body gap are each \qty{0.20}{\milli\metre}.
Single/quad body lengths are \qtylist{20.4;40.6}{\milli\metre}, with pitches
\qtylist{20.6;40.8}{\milli\metre}; neighbouring active edges are separated
by \qty{0.60}{\milli\metre}, and quad chip seams are
\qty{0.20}{\milli\metre}. The selected slim-edge geometry follows the
planar-sensor development motivation.\cite{slim}

Rows have constant pitch within the \qty{550}{\milli\metre} half-length
allocation, with one chip centre at $z=0$ and a common longitudinal phase
between staves. The single-module body is displaced
\qty{0.9}{\milli\metre} tangentially from its active centre.
\Cref{fig:pixel-barrels} shows the layer arrangement and packing.

\begin{table}[!htbp]
\caption{Selected barrel population. Radial offsets are rounded display values;
the printed precision does not specify manufacturing tolerances. Rows means
modules per stave. All dimensions are in \unit{\milli\metre}.}
\label{tab:pixel-barrels}
\centering\small
\begin{tabular}{lrrrrr}
\toprule
Layer & Datum $r$ & Offset $\pm$ & Staves & Rows & Modules\\
\midrule
\csname @@input\endcsname generated/pixel-barrel-rows.tex
\bottomrule
\end{tabular}
\end{table}

\begin{figure}[!htbp]
\centering
\includegraphics[width=.53\linewidth]{pixel-barrels}\par\smallskip
\includegraphics[width=.93\linewidth]{pixel-packing}
\caption{Transverse barrel assembly projection (top) and central single/quad
row packing (bottom). Module bodies overlap tangentially on staggered
radial levels; cable envelopes pass outside the bearing rings. The lower
panel resolves active intervals, guards and gaps about $z=0$.}
\label{fig:pixel-barrels}\label{fig:pixel-packing}
\end{figure}

\subsection{Endcap layout}
\label{sec:pixel-endcaps}\label{sec:pixel-endcap-placement}
The six-ring pattern contains \num{102} singles and \num{50} quads,
covering eight radial chip rows. Smaller modules follow the inner curvature;
quads share support and service interfaces at larger radii. The local axes
are tangential $u$ and radial $v$ around every disc. Module centres account
for their individual mounting face and sensor depth.
\Cref{tab:pixel-endcap-rings} gives the nominal ring parameters.

\begin{table}[!htbp]
\caption{Six-ring source pattern. Radii are the uncorrected radial placement
parameters; actual module centres retain the face/height-dependent transverse
compensation. Each ring occupies both mounting faces and two height levels per
face. Quad rings contain two radial chip rows.}
\label{tab:pixel-endcap-rings}
\centering\small
\begin{tabular}{l l r r r}
\toprule Ring & Family & Nominal $r$ [\unit{\milli\metre}] & Modules & Chips\\\midrule
\csname @@input\endcsname generated/pixel-ring-rows.tex
\bottomrule
\end{tabular}
\end{table}

The station inventory in \cref{tab:pixel-discs} specifies the retained inner
rings and both outer quad rings at each datum. The larger forward active apertures are implemented by retaining the selected
ring populations at each station. Local module positions, face assignments and
sensor depths are preserved within each rigid disc assembly. Disc identifiers
remain stable: disc 10 lies at $|z|=\qty{950}{\milli\metre}$ between discs
2 and 3 and repeats the complete local geometry of disc 2. The table is ordered
by distance from the interaction point. These apertures and datums do not imply
complete acceptance for every trajectory.
Both endcaps contain \num{2616} modules and
\num{5616} chip patches. The negative end uses corresponding signed datums
and support-facing module frames.

\begin{table}[!htbp]
\caption{Selected disc support datums and retained inventory on one side.
Ring numbers refer to the six-ring source pattern; quad rings 5 and 6 remain
at every station. The negative side uses the opposite datum sign and the same
populations. Datums identify rigid assemblies, not manufacturing tolerances.}
\label{tab:pixel-discs}
\centering\small
\begin{tabular}{@{}rrlrrrr@{}}
\toprule Disc & $|z|$ [\unit{\milli\metre}] & Single rings & Singles & Quads & Chips & Circuits\\\midrule
\csname @@input\endcsname generated/pixel-disc-rows.tex
\bottomrule
\end{tabular}
\end{table}

Modules alternate between the two plate faces and between two heights on
each face, keeping overlapping occupied bodies separated in depth. A quad's
four patches share one sensor plane. Sensor centres are at local
$w=\pm4.10$ or $\pm5.75\,\unit{\milli\metre}$; ASICs and graphite contacts
face the plate. An outgoing track crosses the plate before an outward-facing
sensor. The first and last discs are compared in \cref{fig:pixel-endcap}.

\begin{figure}[!htbp]
\centering\includegraphics[width=\linewidth]{pixel-endcap}
\caption{First and last positive discs at a common scale. The first has four
single rings and two quad rings; the last retains single ring 4 and both quad
rings. Colour distinguishes face and height. The views project all sensor
depths onto the transverse plane.}
\label{fig:pixel-endcap}
\end{figure}

\subsection{Supports and cooling}
\label{sec:pixel-staves}\label{sec:pixel-barrel-mounting}
\label{sec:pixel-disc-support}\label{sec:pixel-carrier}
Barrel staves use carbon-foam cores, CFRP skins, graphite spreaders and two
embedded titanium tubes. The \qty{10}{\milli\metre} single and
\qty{24}{\milli\metre} quad spines support wider module-facing plates,
bringing heat from the chip area toward the cooled core.
The support lies radially outside the sensor. Six bearing rings per layer
connect the staves to the barrel cage at $z=\pm110,\pm330,\pm546$
\unit{\milli\metre}, giving \num{24} rings and \num{492} feet.
The station at $z=-\qty{546}{\milli\metre}$ sets the axial datum; the
remaining contacts guide the stave while allowing longitudinal contraction.
Two independent evaporators per stave run in opposite directions between
its ends. Carbon sandwiches and titanium cooling tubes have public pixel
stave precedents.\cite{ibl}

Each endcap plate remains a full annulus from $r=27$ to
$188.5\,\unit{\milli\metre}$, including on discs with fewer active rings.
A per-chip pickup and short graphite stem connect each ASIC to the sandwich
and titanium tube. Graphite orientation defines the thermal-conduction
path.\cite{graphite,tepx} Azimuthal and radial tube routes occupy separate
planes at $w=\pm\qty{1.5}{\milli\metre}$. Each chip row has two independent
half-ring circuits; the station totals are in \cref{tab:pixel-discs}, giving
\num{280} circuits across the two endcaps. Three rim tongues transfer loads
through cone/slot/plane contacts to box rails on a common CFRP carrier.
\Cref{tab:pixel-support,fig:pixel-support} specify the local passive geometry.

\begin{table}[!htbp]
\caption{Support and mounting ledger. Lengths are in \unit{\milli\metre}.
Internal adhesive and inserts replace core material within the stated envelope.}
\label{tab:pixel-support}
\centering\small
\begin{tabular}{@{}p{.35\linewidth}p{.59\linewidth}@{}}\toprule
Component & Selected dimensions\\\midrule
Stave stack, module outward & Interface 0.100; insulation 0.025; graphite 0.200;
CFRP 0.150; foam 4.100; CFRP 0.150; total 4.725. Internal adhesive equivalent 0.080.\\
Barrel bearing rings/feet & Ring axial width 8, radial depth 1; inner clearance 0.5
beyond stave corner. Foot tangential width 4, axial length 8; height follows stagger.\\
Passive stave span & Singles: $z=-551.0$ to $+552.8$; quads: $z=\pm552.8$.\\
Disc sandwich & Core 6.0; each CFRP skin0.15; total 6.3.\\
Per-chip pickup & Area $384\,\unit{\milli\metre\squared}$; graphite 0.30;
cradle 0.075; insulation/interface/adhesive0.025 each; total 0.45.\\
Disc stem and inserts & Stem section $6\times8$; heights 0 or 1.65; quad stems
shift 1 toward module centre. Core insert 4.50; skin window 0.15;
contact area $48\,\unit{\milli\metre\squared}$; adhesive equivalent 0.08.\\
Disc tongues & Width 10; $r=186$--228 at $\phi=90,210,330^{\circ}$.\\
Carrier rails/shell/flanges & Station-following coupling slots; three $8\times8$ box rails, wall 0.4 at
$r=223.7$--231.7; shell $r=231.7$--232.0, $|z|=609$--3136;
end flanges radial 10, axial 2.\\\bottomrule
\end{tabular}
\end{table}

\begin{figure}[!htbp]
\centering\includegraphics[width=.95\linewidth]{pixel-stave}\par\smallskip
\includegraphics[width=.95\linewidth]{pixel-disc-section}
\caption{Stave sections (top) and a radial/axial disc projection (bottom).
Carbon, foam, titanium and coolant form the passive thermal path. The disc
panel superposes both mounting faces and heights; short pickups connect the
chip backs to the plate. Thin layers are specified in the support ledger.}
\label{fig:pixel-support}\label{fig:pixel-stave}\label{fig:pixel-disc-section}
\end{figure}

\begin{table}[!htbp]
\caption{Reference cooling inputs. Diameters/walls are in
\unit{\milli\metre}; flow is per independent tube or circuit.
The common heat allowance is \qty{2.688}{\watt} per chip; the endcap flow
budget uses a factor \num{1.5} stress load.}
\label{tab:pixel-cooling}\centering\small
\begin{tabular}{@{}lrrrr@{}}\toprule
Location & Ti diameter & Ti wall & CO$_2$ [\unit{\degreeCelsius}] & Flow [\unit{\gram\per\second}]\\\midrule
Inner barrel & 2.0 & 0.11 & $-35$ & 1.3\\
Outer barrel & 2.8 & 0.15 & $-35$ & 2.5\\
Endcap, inner chip rows & 2.8 & 0.15 & $-40$ & 1.0\\
Endcap, outer quad chip rows & 2.8 & 0.15 & $-40$ & 1.1\\\bottomrule
\end{tabular}
\end{table}

The barrel tube centres are tangentially at $\pm\qty{2.2}{\milli\metre}$
for inner staves and $\pm\qty{6}{\milli\metre}$ for outer staves. The
\num{164} full-length tubes join \num{82} feed and \num{82} return branches
at each signed end. Thin titanium circuits with evaporative CO$_2$ provide
the reference cooling topology.\cite{cooling,tepx}

\FloatBarrier
\subsection{Services and material representation}
\label{sec:pixel-services}\label{sec:pixel-material}
Barrel electrical services follow each half-stave toward its nearest end;
the central single-chip module joins the positive half. Two serial chains
serve each half-stave, with limits of \num{16} modules and \num{32} chips,
giving \num{164} chains at either end. Cable envelopes grow with the
cumulative module and chain load and pass outside the bearing rings.
Cooling and electrical services collect in twelve sectors in the
$|z|=555$--605\,\unit{\milli\metre} bay, then join the axial reservation
near $r=\qty{212}{\milli\metre}$. Two sectors serve six staves each and
ten serve seven each. \Cref{fig:pixel-barrel-routing} illustrates the local
route; \cref{tab:pixel-routing} records its envelope rules.

\begin{figure}[!htbp]
\centering\includegraphics[width=\linewidth]{pixel-barrel-routing}
\caption{Barrel service routes and cumulative cable footprint. The upper
panel projects the occupied bodies, support, bearings and conceptual pickups.
Counterflow paths are separated radially for display; their physical tubes
are separated tangentially in the core.}
\label{fig:pixel-barrel-routing}
\end{figure}

Endcap tails join ring buses on their own module face. Three radial fans
per face gather them at the disc rim. Electrical chains stop at half-ring
boundaries and contain at most \num{32} chips, independently of the cooling
circuits: the full disc has \num{18} chains, reducing to
\numlist{16;14;12} on the forward stations. A \qty{50}{\milli\metre}-deep
collector carries each load into the axial trunk. The last disc uses an inner
bypass to the rear turn at $|z|=\qty{3150}{\milli\metre}$.
\Cref{fig:pixel-endcap-routing,fig:pixel-service-interfaces} show the face
routes and effective transport cells.

\begin{table}[!htbp]
\caption{Electrical and remote-service ledger. Lengths are in
\unit{\milli\metre}, areas in \unit{\milli\metre\squared}.
Cable footprint and inter-cable packing are separate budgets.}
\label{tab:pixel-routing}\centering\small
\begin{tabular}{@{}p{.35\linewidth}p{.59\linewidth}@{}}\toprule
Component & Allocation and route\\\midrule
Barrel cables & One data and one command link per module, footprint 1 each;
ancillary footprint 35 per chain from its first pickup. Sector width 75\% of
stave pitch, packing 50\%; inner radius 0.5 beyond the outer bearing ring.\\
Half-stave terminal loads & Singles: ancillary 70 plus links 52 at $-z$,54 at
$+z$; quads: ancillary 70 plus links 26.\\
Grouped coolant transport & Feed and return branches reserve diameter4;
a grouped leg for $N$ staves has envelope diameter $4\sqrt{N}$.\\
Endcap face flex & Ring buses $8\times0.20$; module tails $8\times0.12$;
three radial fans per face.\\
Disc collector & $r=27$--190; local outward $z=11.7$--61.7 from its datum.\\
Axial trunk & Reservation $r=190$--234; material cells $r=192$--231.7.
Three 90$^{\circ}$ sectors separated by 30$^{\circ}$ mounting sectors centred
at $\phi=90,210,330^{\circ}$. Barrel load enters at $|z|=605$;
disc collectors follow their station datums; the rearmost disc is logical disc 9.\\
Rear interfaces & Last collector ends at $|z|=3131.7$; bypass $r=27$--188
to rear turn $|z|=3150$. Rear cell extends to $|z|=3300$, $r=680$.
Flange necks at $|z|=609$--611 and 3134--3136 have outer radius 222.\\\bottomrule
\end{tabular}
\end{table}

\begin{figure}[!htbp]
\centering\includegraphics[width=.86\linewidth]{pixel-endcap-routing}
\caption{Disc face buses and independent half-ring cooling paths. The
first/last discs have sixteen/ten cooling circuits. Module outlines preserve
their frames; bus lines, fans and inlet/outlet symbols show service topology.}
\label{fig:pixel-endcap-routing}
\end{figure}

\begin{figure}[!htbp]
\centering\includegraphics[width=\linewidth]{pixel-service-interfaces}
\caption{Effective collection, trunk, final-disc bypass and rear-turn cells.
Carrier-flange necks retain their fixed bounds. Collector cells follow the station datums; the trunk and rear cells retain
the reference constituent inventory within their fixed bounds.}
\label{fig:pixel-service-interfaces}
\end{figure}

The simulation separates sensitive silicon from passive module bodies,
supports and services. Explicit sensor substrates, chip patches, ASICs,
graphite contacts, stave spines, tubes, rings and feet preserve material
ordering around the measurements. Effective components conserve their
constituent inventory as summarized in \cref{tab:pixel-material}.

\begin{table}[!htbp]
\caption{Selected representations and the physical quantities they preserve.
The effective components remain identifiable material allocations.}
\label{tab:pixel-material}
\small\begin{tabular}{p{.29\linewidth}p{.62\linewidth}}
\toprule Component & Representation and scope\\\midrule
Disc sandwich & Explicit annulus and disjoint tongues; constituent-normalized
core for inserts, adhesive and closeouts.\\
Cooling & Explicit Ti/CO$_2$ evaporator tori; summed half-ring length preserved.
Radial legs and bends retain effective inventory.\\
Module pickup & Explicit per-chip pickup stacks and short stems; contact windows and
cradle openings preserved through effective volumes.\\
Disc flex and fittings & Local effective annulus beyond module envelopes;
flex, clip, coupling and bend allowances remain separate inventory terms.\\
Collector and trunk & Effective transport cells preserve conductor, insulation,
titanium and coolant volumes through service handoffs.\\
Carrier & Shell, curved box rails and flanges partitioned to avoid counting
intersections twice; coupling pockets remain reserved.\\\bottomrule
\end{tabular}
\end{table}

Remote cable footprint uses provisional volume fractions of
\qty{10}{\percent} copper, \qty{40}{\percent} polyimide and
\qty{50}{\percent} air, with a separate \qty{50}{\percent} inter-cable
packing allowance. CO$_2$ inventory uses a reference density of
\qty{1.0964}{\gram\per\centi\metre\cubed}, the saturated liquid density
at \qty{-35}{\degreeCelsius}.\cite{nistco2} The model-inventory mass is
\qty{204.747}{\kilogram}, including \qty{159.749}{\kilogram} in downstream
transport cells. The maintained DD4hep geometry repeats the detailed disc assemblies, including
sensors, passive stacks, cooling, mounting tongues and local collectors. Rail
coupling slots follow all ten stations; carrier and service envelopes remain
fixed. Global trunk, bypass and rear cells retain the reference material
inventory. Their cumulative connectivity and capacity for the selected
populations, and the thermal, hydraulic and mechanical margins, remain
unqualified.

The reconstruction comparison uses a separate effective ACTS material
representation: existing material profiles are transported with their assigned
surfaces and the added disc has a homogeneous donor-disc material layer behind
its sensors. That approximation does not replace the detailed DD4hep assembly.
A new Geant4 material map and candidate-specific response validation remain
separate requirements.

\clearpage
% Public external references only. Internal provenance is retained outside the prose.
\begingroup
\renewcommand{\bibname}{Pixel detector references}
\renewcommand{\chapter}[2][]{\subsection*{Pixel detector references}}
\small\raggedright
\begin{thebibliography}{9}
\bibitem{rd53} Stefano Esposito, on behalf of the RD53 Collaboration,
\emph{RD53: Lessons Learned --- A verification perspective},
2023-10-23, slide~5: device generations, ATLAS matrix and approximate die
outline. \href{https://indico.cern.ch/event/1335413/contributions/5621748/attachments/2731411/4764279/RD53\%20Lessons\%20Learned\%20in\%20Verification.pdf}{Public presentation}.

\bibitem{slim} Philipp Weigell, on behalf of the ATLAS Planar Pixel Sensors R\&D
Project, \emph{Recent Results of the ATLAS Upgrade Planar Pixel Sensors R\&D
Project}, arXiv:1210.7661v2 (2013). Section~3.1.1, PDF p.~2: n-in-p junction,
guard and edge-insulation interfaces; \S3.2.1, PDF p.~4: slim-edge sensor
developments. \href{https://cds.cern.ch/record/1491076/files/1210.7661.pdf}{Public paper}.

\bibitem{ibl} ATLAS IBL Collaboration, \emph{Production and Integration of the
ATLAS Insertable B-Layer}, JINST~13 (2018) T05008, arXiv:1803.00844v3.
Section~4.1, Tables~9--10, PDF pp.~38--39: carbon supports and their constituents;
\S5.2.1, Fig.~33, PDF pp.~49--50: titanium cooling;
\S7.3, PDF p.~75: ring/foot mounting interfaces.
\href{https://doi.org/10.1088/1748-0221/13/05/T05008}{Published article};
\href{https://arxiv.org/abs/1803.00844}{public preprint}.

\bibitem{graphite} Richard J.~Lemak, Robert J.~Moskaitis, Yoshiro Nogami and
Katsuhiro Chikuba, \emph{High Performance Pyrolytic Graphite Composite Heat
Spreaders}, IMAPS Advanced Technology Workshop on Thermal Management (2011),
slide~6: in-plane and through-plane thermal conductivity.
\href{https://www.mineralstech.com/docs/default-source/refractories-documents/pyrogenics/pyroid/high-performance-pyrolytic-graphite-composite-heat-spreaders.pdf?sfvrsn=50359b20_2}{Public presentation}.

\bibitem{cooling} Sonia Carr\`a, on behalf of the ATLAS ITk group,
\emph{ATLAS ITk Pixel Outer Endcap CO$_2$ cooling system prototypes},
ATL-UPGRADE-PROC-2024-001, 2024-06-20. Section~1, PDF p.~1: temperature targets;
\S2--4 and Figs.~1,~3, PDF p.~2: thin titanium circuits, joints and prototype
cooling arrangements.
\href{https://cds.cern.ch/record/2901596/files/ATL-UPGRADE-PROC-2024-001.pdf}{Public preprint}.

\bibitem{tepx} V.~Lukashenko, on behalf of the CMS Tracker group,
\emph{Results of CMS Phase-2 TEPX system tests}, JINST~21 (2026) C02027.
Section~1 and Fig.~1, PDF p.~2: composite/titanium supports;
\S5, PDF pp.~5--6: module contact force and thermal interface measurements.
\href{https://doi.org/10.1088/1748-0221/21/02/C02027}{Published article}.

\bibitem{nistco2} NIST Chemistry WebBook, Standard Reference Database~69,
\emph{Carbon dioxide saturation properties}, CAS~124-38-9. Saturation table
from $-40$ to $-30\,\unit{\degreeCelsius}$ in \qty{5}{\degreeCelsius} steps,
$-35\,\unit{\degreeCelsius}$ row, liquid-density column; accessed 2026-09-30. \href{https://webbook.nist.gov/cgi/fluid.cgi?Action=Data&Wide=on&ID=C124389&Type=SatP&Digits=5&THigh=-30&TLow=-40&TInc=5&RefState=DEF&TUnit=C&PUnit=bar&DUnit=kg\%2Fm3&HUnit=kJ\%2Fkg&WUnit=m\%2Fs&VisUnit=uPa*s&STUnit=N\%2Fm}{Public tabulated data}.
\end{thebibliography}
\endgroup

